\chapter{State of the Art}
\label{ch:sota}

% Allow a little extra line stretch so long technical compounds rebreak
% cleanly instead of spilling into the margin (no template packages added).
\setlength{\emergencystretch}{3em}

The problem addressed by this dissertation --- deciding, after a remediation
cycle, whether a previously reported vulnerability is still exploitable --- sits at
the meeting point of two bodies of work that have grown steadily closer over the
past few years: the practice of penetration testing and vulnerability management,
and the automation of security assessment, increasingly driven by artificial
intelligence. These are not separate worlds. A sustained research and industry
effort has been pushing automation ever deeper into the security lifecycle, and the
line between a manual audit and an automated one is dissolving quickly: the very
period that produced autonomous offensive agents also produced defensive
counterparts such as Google DeepMind's \emph{CodeMender}, an AI agent that patches
newly discovered vulnerabilities and proactively rewrites existing source code to
remove whole classes of weakness, with every patch reviewed by a human before it is
upstreamed \cite{codemender}. Against this fast-moving and increasingly crowded
backdrop, the report-driven revalidation of a specific reported finding --- re-running
its own documented reproduction steps rather than a generic attack library --- is one of
the tasks this wave of automation has so far largely passed over, and it is that
omission this chapter builds toward. It first frames the manual, repetitive nature of the
\emph{revalidation} (retest) phase within the vulnerability management lifecycle
(Section~\ref{sec:pentest-lifecycle}). It then surveys the successive waves of
automation applied to security testing, from signature-based scanners to
learning-based autonomous agents (Section~\ref{sec:automation}). The following
two sections concentrate on the technology that makes this work possible: large
language models (LLMs) and the tool-using agents built on top of them
(Section~\ref{sec:llm-agents}), and the extraction of structured information from
the unstructured prose of security reports (Section~\ref{sec:ie}). The chapter
closes with a synthesis that positions the proposed system against related work
and identifies the narrow gap it targets (Section~\ref{sec:gap}).

\section{Penetration testing and the vulnerability management lifecycle}
\label{sec:pentest-lifecycle}

A penetration test is an authorised, goal-directed simulation of an attack whose
purpose is to discover and demonstrate exploitable weaknesses in a system before a
real adversary does \cite{nist80015}. Its output is a report: a structured
account of each finding, its severity, its impact, and the concrete steps needed
to reproduce it. That report is not an end in itself but the entry point to the
\emph{vulnerability management lifecycle} --- the cyclical process of
identifying, prioritising, remediating and \emph{re-verifying}
weaknesses. This section reviews the methodologies that standardise the practice,
isolates the revalidation phase that this dissertation targets, and summarises the
taxonomies used to classify and score the findings that flow through the cycle.

\subsection{Methodologies and standards}
\label{subsec:methodologies}

The discipline has converged on several public methodologies that prescribe how a
test should be scoped, executed and reported. The Penetration Testing Execution
Standard (PTES) \cite{ptes} defines seven phases, from pre-engagement interactions
to reporting, and is deliberately technology-agnostic. The Open Source Security
Testing Methodology Manual (OSSTMM) \cite{osstmm} emphasises measurable, repeatable
verification and introduces the notion of an operational security metric. For web
applications --- the scope of this work --- the OWASP Web Security Testing Guide
(WSTG) \cite{owasp_wstg} provides a catalogue of concrete test cases organised by
vulnerability class. At the institutional level, the NIST Technical Guide to
Information Security Testing and Assessment (SP~800-115) \cite{nist80015} frames
testing within a broader assessment programme and explicitly includes
\emph{re-testing} of remediated findings as part of the assessment rather than an
optional afterthought.

A common thread runs through all four documents: each treats the test as a
\emph{process} that produces evidence, and each assumes that remediation will be
followed by verification. Table~\ref{tab:standards} summarises how the principal
methodologies position the revalidation step. The recurring observation is that,
although every standard expects re-verification, none prescribes \emph{how} it
should be automated --- the activity is described as a human responsibility, to be
repeated in full whenever a fix is deployed.

\begin{table}[htbp]
\footnotesize
\renewcommand{\arraystretch}{1.3}
\begin{center}
\begin{tabular}{|L{3.2cm}|L{4.2cm}|L{5.6cm}|}
\hline
\rowcolor{tema!10}
\bft{Methodology / standard} & \bft{Primary focus} & \bft{Treatment of the retest / revalidation phase} \\ \hline
PTES & End-to-end engagement process (7 phases) & Reporting phase expects verification of fixes; procedure left to the tester. \\ \hline
OSSTMM & Measurable, repeatable verification & Repeatability is a core principle, but retest is a manual re-run of the same checks. \\ \hline
OWASP WSTG & Web-application test cases by class & Provides the concrete checks to re-run; no automation of the re-run itself. \\ \hline
NIST SP~800-115 & Institutional assessment programme & Recommends re-testing of remediated findings as part of the assessment. \\ \hline
\end{tabular}
\end{center}
\caption[Treatment of revalidation in penetration-testing standards]{How the
principal penetration-testing methodologies position the revalidation (retest)
phase. Every standard expects re-verification, yet none specifies how to
automate it.}
\label{tab:standards}
\renewcommand{\arraystretch}{1}
\end{table}

\subsection{Remediation and the revalidation phase}
\label{subsec:revalidation}

Once a report is delivered, the responsibility passes to the development or
operations team, which applies fixes and, in turn, requests confirmation that the
reported weaknesses are genuinely closed. This confirmation --- the
\emph{revalidation} or \emph{retest} --- is the phase this dissertation
automates, and Figure~\ref{fig:lifecycle} situates it within the wider
vulnerability-management cycle. It is uniquely costly for three reasons. First, it is
\emph{repetitive}: each fix triggers a re-run of the original reproduction steps,
often verbatim, against a changed system. Second, it is \emph{judgement-laden}:
deciding that a finding is fixed requires distinguishing a genuine remediation
from an incidental change --- a moved endpoint, an altered response, a disabled
feature --- that merely hides the symptom. Third, it is \emph{poorly supported by
tooling}: the report is prose, the reproduction steps are written for a human, and
no artefact links a finding to a machine-checkable verification procedure. That such
verification is necessary rather than a formality is borne out by remediation data:
across the more than five thousand penetration tests Cobalt runs each year, only
48\% of findings --- and 69\% of the most serious ones --- are ever resolved, at a
median of 67 days \cite{cobalt_pentesting_2025}; a reported fix cannot be assumed to
hold, it must be checked. Compliance regimes compound the pressure. The Payment Card Industry Data Security
Standard, for example, requires that exploitable vulnerabilities found during
penetration testing be corrected and the correction verified by repeat testing
\cite{pcidss}, turning revalidation into a recurring, deadline-bound obligation
rather than a one-off task.

\begin{figure}[htbp]
\centering
\begin{tikzpicture}[
  >={Stealth[length=2.5mm]},
  phase/.style={draw=gris50, thick, rounded corners, align=center, minimum width=3.1cm, minimum height=1.1cm, fill=gris95, font=\small},
  hl/.style={phase, draw=rojo, very thick, fill=rojo!8},
]
\node[phase] (id)  at (0,2.3) {\bft{Identify}\\\footnotesize penetration test};
\node[phase] (pr)  at (6,2.3) {\bft{Prioritise}\\\footnotesize severity (CVSS)};
\node[phase] (rem) at (6,0)   {\bft{Remediate}\\\footnotesize apply fixes};
\node[hl]    (rev) at (0,0)   {\bft{Revalidate}\\\footnotesize retest};
\draw[->, gris50, thick]    (id)  -- (pr);
\draw[->, gris50, thick]    (pr)  -- (rem);
\draw[->, gris50, thick]    (rem) -- (rev);
\draw[->, rojo, very thick] (rev) -- (id);
\node[font=\footnotesize\itshape, text=rojo, below=4pt of rev] {automated in this work};
\end{tikzpicture}
\caption[The vulnerability-management lifecycle]{The vulnerability-management
lifecycle: a finding moves from identification through prioritisation and
remediation to revalidation --- the retest that confirms the fix before the cycle
repeats. The revalidation phase (highlighted) is the one this work automates.}
\label{fig:lifecycle}
\end{figure}

The consequence is a verification bottleneck. The creative, exploratory part of an
audit (finding the vulnerability) is performed once, whereas the mechanical
part --- proving it is gone --- is repeated on every remediation cycle, yet the two
are supported by the same manual workflow. It is precisely this asymmetry that
motivates applying automation, and specifically language-model-driven automation,
to the revalidation phase.

\subsection{Vulnerability taxonomies and severity scoring}
\label{subsec:taxonomies}

Findings that flow through the lifecycle are described using a small set of widely
adopted reference frameworks, and any tool that ingests reports must be able to
speak their language. Individual issues are catalogued as Common Vulnerabilities
and Exposures (CVE) identifiers, while the underlying weakness classes
are organised by the Common Weakness Enumeration (CWE). Severity is
communicated with the Common Vulnerability Scoring System (CVSS) \cite{cvss},
whose numeric vectors drive remediation priority and, indirectly, which findings
are retested first. For web applications, the OWASP Top~10
provides a consensus ranking of the most critical risk categories --- injection,
broken access control, and related classes --- that dominate real reports.
Finally, the MITRE ATT\&CK knowledge base \cite{mitre_attack} catalogues adversary
tactics and techniques observed in the wild, offering a common vocabulary that
links a reported finding to the attacker behaviour it enables. These frameworks
are relevant here because a revalidation system must map free-text findings onto
them to reason about what a finding \emph{is} and how it should be re-verified.

\section{Automation of security testing}
\label{sec:automation}

Efforts to automate security assessment have proceeded in waves, each trading a
different balance of coverage, precision and autonomy. This section reviews three:
scanners that automate detection but not reasoning, breach-and-attack-simulation
platforms that automate the execution of known adversary behaviour, and
learning-based systems that attempt to automate the planning of an attack itself.
Figure~\ref{fig:waves} summarises these three waves and the trend toward greater
autonomy they trace.

\begin{figure}[htbp]
\centering
\begin{tikzpicture}[
  >={Stealth[length=3mm]},
  wave/.style={draw=gris50, thick, rounded corners, align=center, text width=3.5cm, minimum height=2.4cm, fill=gris95, font=\footnotesize, inner sep=5pt},
]
\node[wave] (w1) at (0,0)   {\bft{Wave 1}\\[2pt] Vulnerability scanners / DAST\\[4pt] \emph{match known patterns}\\[2pt] Nessus, OpenVAS, ZAP};
\node[wave] (w2) at (4.7,0) {\bft{Wave 2}\\[2pt] Breach \& attack simulation\\[4pt] \emph{replay known techniques}\\[2pt] CALDERA, Atomic Red Team};
\node[wave] (w3) at (9.4,0) {\bft{Wave 3}\\[2pt] Autonomous / LLM agents\\[4pt] \emph{plan, act and adapt}\\[2pt] PentestGPT, XBOW, \emph{Mythos}};
\draw[->, gris50, thick] (w1) -- (w2);
\draw[->, gris50, thick] (w2) -- (w3);
\draw[->, thick, tema] ($(w1.south west)+(0,-0.6)$) -- node[below, font=\footnotesize\itshape, text=gris25] {increasing autonomy, decreasing human involvement} ($(w3.south east)+(0,-0.6)$);
\end{tikzpicture}
\caption[The three waves of security-testing automation]{The three waves of
security-testing automation, from pattern-matching scanners to autonomous
language-model agents. Each wave delegates more of the task to the machine; the
system proposed in this work deliberately reintroduces a human decision on every
action.}
\label{fig:waves}
\end{figure}

\subsection{Vulnerability scanners and dynamic analysis}
\label{subsec:scanners}

The most widespread form of automation is the vulnerability scanner. Network and
host scanners such as Nessus and OpenVAS match
observed service banners and configurations against databases of known issues,
whereas dynamic application-security-testing (DAST) tools, among them OWASP
ZAP and Burp Suite, crawl a running web application
and inject payloads to elicit signs of injection, cross-site scripting or
misconfiguration. These tools are fast and
broad, but their limitations are well documented. A controlled evaluation of
eleven black-box web scanners found that entire vulnerability classes --- notably
stored cross-site scripting and logic flaws --- are systematically missed, and
that results vary widely across tools on the same target \cite{doupe2010} --- a
substantial fraction of real vulnerabilities is left undetected. Two properties matter for revalidation in particular:
scanners reason at the level of \emph{patterns}, not of a specific reported
finding with its own reproduction steps; and they are prone to false positives,
which is unacceptable when the output is a binary ``fixed / still-open'' verdict
that gates a compliance decision. A scanner can suggest that a class of problem
may exist; it cannot, on its own, revalidate that \emph{this} finding, described
in \emph{this} report, is closed.

\subsection{Breach and attack simulation}
\label{subsec:bas}

A second wave automates not detection but the \emph{execution} of known adversary
behaviour. Breach-and-attack-simulation (BAS) platforms encode attacker techniques
as reusable, executable actions and replay them against a target to test whether
defences hold. MITRE CALDERA orchestrates such actions as an
autonomous red-team agent whose behaviours are mapped to ATT\&CK techniques, while
Atomic Red Team provides a library of small, technique-scoped
tests intended to be run repeatably. The intellectual foundation for this
approach was laid by work on automated adversary emulation, which showed that a
planner could chain individual techniques into a coherent attack sequence
\cite{applebaum2016}. BAS shares one important trait with the present work --- the
idea that verification should be an executable, repeatable procedure rather than a
manual re-run --- but it differs in one important respect: its action library is
authored in advance and is agnostic to any particular report. A BAS platform
answers ``can technique~T succeed here?''; revalidation must answer ``is the
specific finding~F, as reported, still exploitable?'', a question tied to the
reproduction steps of an individual report rather than to a catalogue of generic
techniques.

\subsection{Autonomous and learning-based penetration testing}
\label{subsec:autonomous}

A third line of research treats penetration testing as a sequential
decision-making problem and applies automated planning and reinforcement learning
to it. Early work modelled the attacker's uncertainty explicitly, casting
penetration testing as a partially observable Markov decision process so that a
planner could reason about what it did not yet know \cite{sarraute2012}. The line
that followed it broadened into a spectrum of simulated-pentesting models of
increasing fidelity --- from classical planning to game-theoretic formulations ---
and, more recently, into reinforcement-learning agents trained to select
exploitation actions directly from the state of a simulated network. These
approaches demonstrate that the
\emph{planning} of an attack can be automated, but they operate almost exclusively
in simulated or abstracted environments, they require a formal model of the target,
and they optimise for \emph{finding} a path to compromise rather than for
\emph{confirming} the status of a pre-existing, human-described finding. The gap
between a formally modelled simulation and the messy reality of re-verifying a
report written in natural language is exactly where language models change the
picture.

\section{Large language models and autonomous agents}
\label{sec:llm-agents}

The final wave of automation is driven by large language models. Their relevance
to this work is twofold: they can interpret the free-text prose of a report, and
they can be turned into agents that reason about and act on a live system. This
section traces the path from language models to tool-using agents, reviews their
recent and rapid application to offensive security, and examines the reliability
concerns that make human oversight indispensable.

\subsection{From language models to tool-using agents}
\label{subsec:llm-foundations}

Modern language models are built on the transformer architecture \cite{vaswani2017},
and their most consequential property emerged with scale: sufficiently large models
perform new tasks from a natural-language description and a few examples, without
task-specific training \cite{brown2020}. Two developments turned this capability
into a basis for autonomous action. Prompting a model to produce intermediate
reasoning steps was shown to improve performance on multi-step problems, and --- most
directly relevant here --- the ReAct paradigm interleaved those reasoning traces with
\emph{actions}, letting a model call external tools, observe their results and
adjust its plan accordingly \cite{yao2023react}, with related work showing that
models can also learn \emph{when} and \emph{how} to invoke a tool. Together these results define the tool-using agent: a
loop in which a model proposes an action, an external environment executes it, and
the observed output conditions the next proposal. This loop is the computational
pattern on which the present system is built, with the deliberate addition of a
human approval gate at each step.

\subsection{Language-model agents for offensive security}
\label{subsec:llm-offensive}

The application of tool-using agents to offensive security is recent and moving
quickly, and such systems have grown rapidly across the security lifecycle.
PentestGPT decomposes a penetration test
into interacting reasoning, generation and parsing modules and shows that an LLM
can guide the earlier stages of a test, while also documenting where the model
loses track of context on longer engagements \cite{deng2024pentestgpt}. In parallel,
an approach explicitly framed around a human-in-the-loop workflow paired a
language model with a security tester to assist with penetration-testing tasks,
arguing that human oversight is essential when the model acts on a real system
\cite{happe2023}. Stronger results followed: autonomous
agents were shown to exploit web vulnerabilities with no prior knowledge of them
and, subsequently, to exploit one-day vulnerabilities from their public descriptions
alone \cite{fang2024oneday}; others extended the idea to broader libraries of
post-compromise attacker actions.

These capabilities have moved from research prototypes into operational tools that
work in much the same way as the retest agent developed here: an agent proposes an
action, runs it, observes the result and iterates. The offensive systems introduced
in Section~\ref{sec:motivation} are instances of this shift --- the autonomous
penetration tester XBOW \cite{xbow} and the frontier discovery-and-exploitation
models exemplified by Mythos \cite{anthropic_mythos, aisi_mythos}. One detail of XBOW is worth carrying
forward: it pairs its language-model-driven exploration with \emph{deterministic}
validation of each finding and human review before disclosure, a discipline close in
spirit to the one adopted here. What all these systems share, however, is their
target: the \emph{discovery} or \emph{exploitation} of vulnerabilities. Revalidation
is adjacent but distinct --- it re-runs the reproduction steps of an
\emph{already-reported} finding to decide whether a fix holds --- and it can
therefore afford, and this work argues it should require, the human gate on every
action that the autonomy-maximising offensive tools deliberately relax.

Two lessons from this literature shape the present work. The capability is real and offensively useful, which makes
\emph{unsupervised} execution against real systems irresponsible; and the same
capability, when constrained, is exactly what is needed to re-execute a finding's
reproduction steps and observe whether they still succeed.

\subsection{Reliability and safety of language-model agents}
\label{subsec:llm-safety}

Deploying an autonomous agent against a live target raises reliability and safety
concerns that directly motivate this system's design. Language models are prone to
\emph{hallucination} --- producing fluent, confident output that is unsupported by
evidence \cite{ji2023hallucination}. In a revalidation context this failure is
acute: a model that rationalises an inconclusive result into a confident ``fixed''
verdict is actively harmful, because it retires a still-open vulnerability. Agents
that ingest untrusted content are further exposed to indirect prompt injection,
where data returned by a tool subverts the model's instructions
\cite{greshake2023} --- a serious risk when the ``data'' is the response of a
system under test. These concerns argue for three design commitments adopted in
this work and revisited in later chapters: a \emph{human-in-the-loop} gate that
approves every action before it executes; \emph{sandboxing} that confines the agent
to the authorised target; and \emph{structured, validated output} together with a
conservative bias toward an \emph{inconclusive} verdict whenever the evidence does
not unambiguously support a conclusion.

\section{Information extraction from security reports}
\label{sec:ie}

Before any retest can run, the system must convert a report --- typically a PDF
written for human readers --- into structured findings with explicit reproduction
steps. This is an information-extraction problem, and it has an established research
lineage in the adjacent domain of cyber threat intelligence (CTI). Earlier systems
extracted attacker tactics and techniques from unstructured threat reports and
mapped them onto formal representations, or retrieved ATT\&CK techniques from such
reports by supervised classification; EXTRACTOR went further and distilled attack
behaviour from prose into provenance-style graphs \cite{satvat2021extractor}. These systems
established that security-relevant structure can be recovered from free text, but
they relied on domain-specific pipelines and targeted a fixed output schema. Large
language models change the economics of this task: given only a description of the
target structure, a single model can extract a rich, nested schema from
heterogeneous report layouts. The approach adopted in this dissertation depends on
that extraction being \emph{schema-constrained} --- techniques such as grammar- and
JSON-schema-guided decoding constrain a model's output to conform to a target
structure \cite{willard2023} --- and validated, so that malformed or hallucinated
output is rejected before it reaches persistence rather than silently corrupting
every downstream verdict.

\section{Synthesis and research gap}
\label{sec:gap}

The preceding sections describe a field with many adjacent capabilities but a
gap at its centre. Standards mandate revalidation without automating it
(Section~\ref{subsec:revalidation}). Scanners automate detection but reason about
patterns, not about a specific reported finding, and their false positives are
incompatible with a gating verdict (Section~\ref{subsec:scanners}). BAS platforms
automate repeatable execution but from a report-agnostic library
(Section~\ref{subsec:bas}). Autonomous planners automate attack discovery but in
simulated environments and toward compromise rather than confirmation
(Section~\ref{subsec:autonomous}). LLM agents supply the missing ingredient
(interpreting a natural-language finding and acting on a live system), but the
offensive literature pursues \emph{exploitation}, unsupervised, whereas safety
considerations demand the opposite (Section~\ref{subsec:llm-safety}).

Table~\ref{tab:comparison} contrasts these families along the dimensions that
matter \emph{for revalidation specifically}: whether the approach is driven by the
content of an individual report, its execution model, the granularity of human
oversight, and whether it produces a per-finding remediation verdict. These are not
axes of general merit --- each family is engineered for a different goal, and most
were never meant to answer the revalidation question at all. Read that way, the
table reports \emph{fit} rather than quality: no existing family is aimed at
report-driven, human-gated revalidation, which is the narrow slot this work
occupies.

That slot is a deliberate specialisation, and it is bought at a real price.
\emph{revalid} does none of what places the offensive systems of
Section~\ref{subsec:llm-offensive} at the frontier: it does not discover new
vulnerabilities, it does not synthesise exploits, and, because it stops for a human
decision before every command, it neither runs unattended nor scales like a scanner.
On raw autonomous capability it is far behind tools such as XBOW or \emph{Mythos},
and it does not set out to compete with them. Its contribution is narrower and more
concrete: to treat the \emph{revalidation} of an already-reported finding as a
first-class task --- re-executing that finding's reproduction steps, gating every
command on a human decision, and returning a conservative, evidence-linked verdict,
which is the combination the surveyed families are not built to provide. The closest
prior discipline is XBOW's, which already pairs its exploration with deterministic
per-finding validation and human review; it differs in the two respects that define
this niche --- that review is post-hoc, applied before disclosure rather than before
each command, and its object is a newly discovered finding rather than the
reproduction steps of one already reported. The report-driven, human-gated system
that realises this is proposed in Section~\ref{sec:solution}. Whether the
specialisation is worth the autonomy it gives up is a design bet this dissertation
makes deliberately: the chapters that follow develop it and then put it to the test,
in the preliminary evaluation of Chapter~\ref{ch:evaluation}.

\begin{table}[htbp]
\scriptsize
\renewcommand{\arraystretch}{1.4}
\begin{center}
\begin{tabular}{|L{2.4cm}|L{2.7cm}|L{2.2cm}|L{2.2cm}|L{2.6cm}|}
\hline
\rowcolor{tema!10}
\bft{Approach} & \bft{Primary goal} & \bft{Driven by a specific report?} & \bft{Human oversight} & \bft{Per-finding remediation verdict?} \\ \hline
Vulnerability scanners / DAST & Find issue classes across a target & No (pattern-based) & Review of results & Not its aim (flags issue classes) \\ \hline
Breach \& attack simulation & Test whether known techniques succeed & No (fixed library) & Configuration-time & Pass/fail of a technique \\ \hline
Autonomous / RL pentesting & Find a path to compromise & No (formal model) & None (autonomous) & Not its aim (targets compromise) \\ \hline
Offensive LLM agents & Discover and exploit vulnerabilities & Partially & Post-hoc review, not per action & Not its aim (targets exploitation) \\ \hline
\emph{revalid} (this work) & Re-verify a reported fix & Yes & Per-command gate & Yes: open / fixed / inconclusive \\ \hline
\end{tabular}
\end{center}
\caption[Automated security-testing families, contrasted on the axes revalidation
needs]{The families surveyed in this chapter, compared only along the axes a
revalidation tool must satisfy. Each pursues a different primary goal, so the table
measures \emph{fit} to the revalidation task, not general capability:
\emph{revalid} is the only entry built for report-driven, human-gated
re-verification because that is the only thing it does. In exchange it forgoes the
discovery, exploitation and unattended scale that place the other families --- the
offensive LLM agents above all --- well ahead of it on raw capability.}
\label{tab:comparison}
\renewcommand{\arraystretch}{1}
\end{table}
